% ============================================================
% บทที่ 3: เกตเวย์โครงข่ายและการบูรณาการระบบอัตโนมัติ Zigbee2MQTT
% ============================================================
\chapter[เกตเวย์โครงข่ายและการบูรณาการระบบอัตโนมัติ Zigbee2MQTT]{เกตเวย์โครงข่าย\\และการบูรณาการระบบอัตโนมัติ Zigbee2MQTT}
\label{chap:ch03}

% ------------------------------------------------------------
% แผนบริหารการสอนประจำบทที่ 3 (หน้า 1/2)
% ------------------------------------------------------------
\section*{แผนบริหารการสอนประจำบทที่ 3}
\addcontentsline{toc}{section}{แผนบริหารการสอนประจำบทที่ 3}

\begin{planbox}{หัวข้อเนื้อหาประจำบท}
\begin{enumerate}[topsep=1pt, itemsep=0.5pt, partopsep=0pt, parsep=0pt, leftmargin=1.5em]
    \item สถาปัตยกรรมตัวประสานงานโครงข่ายฮาร์ดแวร์ CC2652P และ EFR32MG21
    \item สถาปัตยกรรมซอฟต์แวร์สะพานข้อมูล Zigbee2MQTT และการบริหารจัดการเครือข่าย
    \item โครงสร้างหัวข้อข้อมูล MQTT และข้อกำหนดโครงสร้างเพย์โหลดแบบ JSON Telemetry
    \item การเชื่อมโยงข้อมูลสู่แพลตฟอร์มบริหารจัดการฟาร์มอัจฉริยะ Home Assistant และ Node-RED
    \item ความปลอดภัยของเครือข่าย กุญแจรหัสผ่าน และการป้องกันการแทรกแซงสัญญาณ
\end{enumerate}
\end{planbox}

\begin{planbox}{วัตถุประสงค์เชิงพฤติกรรม}
\begin{enumerate}[topsep=1pt, itemsep=0.5pt, partopsep=0pt, parsep=0pt, leftmargin=1.5em]
    \item ติดตั้งและตั้งค่าเฟิร์มแวร์ตัวประสานงาน CC2652P ได้อย่างถูกต้อง
    \item กำหนดค่าคอนฟิกูเรชันของโปรแกรม Zigbee2MQTT ให้สื่อสารกับโหนดในฟาร์มได้
    \item ออกแบบผังโครงสร้างหัวข้อ MQTT สำหรับข้อมูลโทรมาตรและคำสั่งควบคุมได้อย่างมีมาตรฐาน
    \item สร้างระบบอัตโนมัติเปิดปิดวาล์วน้ำตามเงื่อนไขของเซนเซอร์ใน Home Assistant ได้
    \item กำหนดค่าความปลอดภัยและการเข้ารหัส AES-128 บนเครือข่าย Zigbee ได้อย่างรัดกุม
\end{enumerate}
\end{planbox}

\newpage

% ------------------------------------------------------------
% แผนบริหารการสอนประจำบทที่ 3 (หน้า 2/2)
% ------------------------------------------------------------
\begin{planbox}{วิธีสอนและกิจกรรมการเรียนการสอน}
\begin{enumerate}[topsep=1pt, itemsep=0.5pt, partopsep=0pt, parsep=0pt, leftmargin=1.5em]
    \item การบรรยายหลักการทำงานของเกตเวย์และการแปลงโพรโทคอลระดับบริดจ์
    \item การสาธิตการแฟลชเฟิร์มแวร์ Z-Stack และการตั้งค่า Zigbee2MQTT บน Raspberry Pi
    \item การทดลองจับคู่โหนดเซนเซอร์ใหม่เข้าสู่โครงข่ายและการตรวจสอบแดชบอร์ดผังเมช
    \item การฝึกปฏิบัติการเขียน Automation Script สั่งงานรีเลย์ควบคุมน้ำใน Node-RED
    \item การทดสอบจำลองเหตุการณ์ขาดการเชื่อมต่ออินเทอร์เน็ตเพื่อประเมินเสถียรภาพระดับโลคอล
\end{enumerate}
\end{planbox}

\begin{planbox}{สื่อการเรียนการสอน}
\begin{enumerate}[topsep=1pt, itemsep=0.5pt, partopsep=0pt, parsep=0pt, leftmargin=1.5em]
    \item บอร์ดคอมพิวเตอร์แม่ข่าย Raspberry Pi 4 หรือ Single Board Computer
    \item ตัวประสานงาน USB Zigbee 3.0 Coordinator (ชิปเซต TI CC2652P) พร้อมเสา 5dBi
    \item ซอฟต์แวร์ Zigbee2MQTT, Eclipse Mosquitto MQTT Broker และ Home Assistant OS
    \item โปรแกรมเครื่องมือทดสอบ MQTT Explorer สำหรับตรวจสอบข้อความ JSON
    \item ผังโครงสร้างเครือข่ายและสเปกการตั้งค่าพารามิเตอร์ด้านความปลอดภัย
\end{enumerate}
\end{planbox}

\begin{planbox}{การวัดและประเมินผล}
\begin{enumerate}[topsep=1pt, itemsep=0.5pt, partopsep=0pt, parsep=0pt, leftmargin=1.5em]
    \item การประเมินผลการติดตั้งและตั้งค่าเกตเวย์ให้สามารถจับคู่อุปกรณ์ได้อย่างสมบูรณ์
    \item การตรวจความถูกต้องของโครงสร้าง JSON Telemetry Payload ที่ส่งขึ้น MQTT Broker
    \item การประเมินผลการทำงานของระบบอัตโนมัติในการสั่งงานอุปกรณ์ขับโหลด
    \item การทดสอบความเข้าใจเกี่ยวกับกลไกความปลอดภัยและการเข้ารหัสข้อมูลในโครงข่าย
\end{enumerate}
\end{planbox}

\clearpage

% ------------------------------------------------------------
% เนื้อหาบทที่ 3
% ------------------------------------------------------------
\section{สถาปัตยกรรมฮาร์ดแวร์ตัวประสานงานเครือข่าย}

หัวใจสำคัญของการควบคุมและบริหารจัดการโครงข่าย Zigbee คือ \textbf{ตัวประสานงานโครงข่าย (Zigbee Coordinator)}\index{Coordinator@Coordinator (ตัวประสานงาน)} ซึ่งทำหน้าที่เป็นสะพานเชื่อมต่อระหว่างคลื่นวิทยุ IEEE 802.15.4 เข้ากับระบบเครือข่ายคอมพิวเตอร์ไอพี (Ethernet / Wi-Fi) ชิปเซตที่ได้รับการยอมรับในระดับอุตสาหกรรมสูงสุดในปัจจุบัน ได้แก่

\begin{protocolbox}[การเลือกชิปเซตสำหรับเกตเวย์เกษตรกรรม]
คุณลักษณะเชิงเปรียบเทียบของชิปเซต Coordinator
\begin{enumerate}
    \item \textbf{Texas Instruments CC2652P}\index{CC2652P@CC2652P (ชิปเซตตัวประสานงาน)} ไมโครคอนโทรลเลอร์แกน ARM Cortex-M4F พร้อมภาคขยายกำลังส่งในตัว (Power Amplifier หรือ PA) ให้กำลังส่งสูงสุดถึง $+20$ dBm รองรับการเชื่อมต่อโหนดโดยตรงได้มากกว่า 100 โหนด และรองรับโครงข่ายเมชรวมมากกว่า 200 โหนด มีเสถียรภาพสูงสุดกับระบบ Z-Stack
    \item \textbf{Silicon Labs EFR32MG21}\index{EFR32MG21@EFR32MG21} สถาปัตยกรรม ARM Cortex-M33 กำลังส่ง $+20$ dBm มีจุดเด่นด้านประสิทธิภาพการประมวลผลและการรองรับโพรโทคอล Thread และ Matter ในอนาคต
\end{enumerate}
\end{protocolbox}

สำหรับสภาพแวดล้อมฟาร์มเกษตรกรรม แนะนำให้ใช้ตัวประสานงานแบบเชื่อมต่อผ่านสายแลนพร้อมระบบจ่ายไฟผ่านสายสัญญาณ (Power over Ethernet หรือ PoE) มากกว่าแบบต่อตรงผ่านพอร์ต USB เข้ากับคอมพิวเตอร์ เนื่องจากสามารถติดตั้งเกตเวย์ไว้ในตำแหน่งเสาสูงกึ่งกลางแปลงปลูกเพื่อรับสัญญาณรอบทิศทางได้ดีที่สุด และลดสัญญาณรบกวนความถี่สูงจากพอร์ต USB 3.0

\begin{figure}[H]
    \centering
    \includegraphics[width=0.90\textwidth]{figures/zigbee_greenhouse_ghibli.jpg}
    \caption{ระบบควบคุมอัตโนมัติภายในโรงเรือนเกษตรกรรมอัจฉริยะที่บูรณาการเกตเวย์ Zigbee และเซนเซอร์สิ่งแวดล้อม}
    \label{fig:zigbee_greenhouse}
\end{figure}

\section{สถาปัตยกรรมซอฟต์แวร์ Zigbee2MQTT}

โปรแกรม \textbf{ซิกบีทูเอ็มคิวทีที (Zigbee2MQTT)}\index{Zigbee2MQTT@Zigbee2MQTT} เป็นสะพานซอฟต์แวร์โอเพนซอร์สที่ทำหน้าที่แปลชุดข้อมูลและคำสั่งจากเฟรม ZCL (Zigbee Cluster Library)\index{ZCL@ZCL (Zigbee Cluster Library)} บนคลื่นวิทยุ ให้กลายเป็นข้อความ JSON บนโพรโทคอล MQTT\index{MQTT@MQTT} โดยตรง ทำให้ระบบไอโอทีสามารถสื่อสารกับอุปกรณ์ Zigbee จากผู้ผลิตต่างยี่ห้อกันได้มากกว่า 3,000 รุ่นอย่างไร้รอยต่อ

\begin{codebox}[ตัวอย่างไฟล์คอนฟิกูเรชัน configuration.yaml ของ Zigbee2MQTT]{yaml}
# Zigbee2MQTT configuration for smart farm
homeassistant: true
permit_join: false # Disabled after pairing for security

mqtt:
  base_topic: cmu_aiot/zigbee
  server: mqtt://10.0.0.1:1883
  user: farm_admin
  password: secure_farm_password_2027

serial:
  port: /dev/ttyUSB0
  adapter: zstack
  baudrate: 115200

advanced:
  pan_id: 0x1a62
  channel: 25 # Channel 25 avoids Wi-Fi channel 1, 6, 11 interference
  network_key:
    - 0x01
    - 0x03
    - 0x05
    - 0x07
    - 0x09
    - 0x0B
    - 0x0D
    - 0x0F
    - 0x02
    - 0x04
    - 0x06
    - 0x08
    - 0x0A
    - 0x0C
    - 0x0E
    - 0x10
\end{codebox}

\section{ข้อกำหนดหัวข้อและโครงสร้างข้อมูล JSON Telemetry}

การกำหนดโครงสร้าง \textbf{หัวข้อข้อมูล (MQTT Topics)} และเพย์โหลด ต้องสอดคล้องกับมาตรฐานการจัดการข้อมูลไอโอที เพื่อให้ระบบแสดงผล แดชบอร์ด และฐานข้อมูล Time-Series สามารถสืบค้นข้อมูลได้อย่างรวดเร็ว

\begin{protocolbox}[มาตรฐานโครงสร้าง MQTT Topic สำหรับฟาร์มเกษตรแม่นยำ]
โครงสร้างหัวข้อแบ่งออกเป็น 3 ระดับหลัก
\begin{enumerate}
    \item \textbf{Telemetry Stream} ส่งข้อมูลพารามิเตอร์เซนเซอร์จากโหนดขึ้นสู่โบรกเกอร์
    \begin{quote}
        \texttt{cmu\_aiot/zigbee/zone1/greenhouseA/soil\_node\_01}
    \end{quote}
    \item \textbf{Command Topic} รับคำสั่งจากระบบอัตโนมัติไปยังอุปกรณ์ขับโหลด (Actuator)
    \begin{quote}
        \texttt{cmu\_aiot/zigbee/zone1/greenhouseA/valve\_01/set}
    \end{quote}
    \item \textbf{Device State Topic} รายงานสถานะปัจจุบันของอุปกรณ์ขับโหลด
    \begin{quote}
        \texttt{cmu\_aiot/zigbee/zone1/greenhouseA/valve\_01}
    \end{quote}
\end{enumerate}
\end{protocolbox}

\begin{codebox}[ตัวอย่างเพย์โหลด JSON Telemetry จากโหนดตรวจวัดสภาพดินและบรรยากาศ]{json}
{
  "device_id": "ZB_SOIL_NODE_07",
  "zone": "zone1",
  "battery_voltage": 3.62,
  "battery_percent": 94,
  "linkquality": 142,
  "soil_moisture": 42.8,
  "soil_temperature": 27.4,
  "soil_ec": 1150,
  "soil_ph": 6.45,
  "air_temperature": 29.8,
  "air_humidity": 68.2,
  "vpd_kpa": 1.34,
  "valve_status": "CLOSED",
  "timestamp": 1790156000
}
\end{codebox}

\section{การบูรณาการระบบอัตโนมัติด้วย Home Assistant และ Node-RED}

เมื่อข้อมูลเซนเซอร์ถูกส่งผ่านเกตเวย์เข้าสู่ระบบ MQTT Broker แพลตฟอร์มบริหารจัดการกลาง \textbf{โฮมแอสซิสแทนท์ (Home Assistant)} จะตรวจพบอุปกรณ์โดยอัตโนมัติ (MQTT Auto Discovery) ผู้ดูแลระบบสามารถเขียนกฎการควบคุมอัตโนมัติ (Automation Rules) เช่น

\begin{agribox}[ตัวอย่างตรรกะควบคุมการให้น้ำอัตโนมัติ]
การควบคุมโซลินอยด์วาล์วให้น้ำอัตโนมัติตามความต้องการของพืช
\begin{enumerate}
    \item \textbf{เงื่อนไขเริ่มต้นการให้น้ำ} หากความชื้นดินลดลงต่ำกว่าร้อยละ 35 และเวลาอยู่ในช่วง 06.00--08.00 น. หรือ 16.00--18.00 น. ให้สั่งเปิดโซลินอยด์วาล์วน้ำ
    \item \textbf{เงื่อนไขหยุดการให้น้ำ} เมื่อความชื้นดินเพิ่มขึ้นถึงร้อยละ 55 หรือระยะเวลาเปิดวาล์วต่อเนื่องครบ 15 นาที ให้สั่งปิดวาล์วทันทีเพื่อความปลอดภัย (Fail-safe Timer)
    \item \textbf{เงื่อนไขระงับการทำงาน} หากเซนเซอร์น้ำฝนตรวจพบฝนตก หรือค่า VPD ต่ำกว่า 0.4 kPa ให้ระงับการให้น้ำเพื่อป้องกันรากพืชเน่า
\end{enumerate}
\end{agribox}

\clearpage

% ------------------------------------------------------------
% คำถามทบทวนท้ายบทที่ 3
% ------------------------------------------------------------
\section*{คำถามทบทวนท้ายบทที่ 3}
\addcontentsline{toc}{section}{คำถามทบทวนท้ายบทที่ 3}

\begin{enumerate}
    \item เพราะเหตุใดการเลือกช่องสัญญาณความถี่ (Channel) ในเครือข่าย Zigbee 3.0 จึงแนะนำให้ใช้ช่องที่ 15, 20 หรือ 25 เมื่อมีการใช้งานโครงข่าย Wi-Fi 2.4 GHz ในพื้นที่เดียวกัน
    \item จงอธิบายบทบาทและประโยชน์ของการใช้โปรแกรม Zigbee2MQTT ในการเชื่อมโยงโหนดเซนเซอร์เข้ากับระบบฐานข้อมูลของฟาร์มเกษตรกรรม
    \item ออกแบบโครงสร้างหัวข้อข้อมูล MQTT (Topic Architecture) สำหรับแปลงผลไม้ขนาด 5 โซน โดยแต่ละโซนประกอบด้วยโหนดเซนเซอร์ดิน 10 ตัว และวาล์วน้ำ 4 ตัว
    \item หากเกตเวย์ขาดการเชื่อมต่ออินเทอร์เน็ตภายนอก ระบบควบคุมอัตโนมัติที่ทำงานบนเซิร์ฟเวอร์โลคอลจะยังคงทำงานได้หรือไม่ เพราะเหตุใด
    \item อธิบายความสำคัญของการเปลี่ยนรหัสกุญแจเครือข่าย (Network Key) เริ่มต้นจากโรงงานในการติดตั้งเครือข่ายไร้สายในแปลงเกษตรขนาดใหญ่
\end{enumerate}

\clearpage
